% ─────────────────────────────────────────────────────────────────────────────
% CHAPTER 7: Conclusion and Future Scope
% ─────────────────────────────────────────────────────────────────────────────

This concluding chapter summarizes the key engineering contributions of the GroDealz project, outlines practical takeaways from deploying the platform, discusses current limitations, and identifies promising avenues for future research and enhancement.

\section{Summary of Contributions}

This project set out to bridge the gap between traditional offline grocery buying and modern quick-commerce by building a software system that empowers both everyday consumers and local neighborhood merchants. The main achievements of this work include:

\begin{enumerate}
    \item \textbf{Autonomous AI Grocery Shopping Agent (Module A):} Designed and implemented a conversational agent on \texttt{/ai-planner} that understands informal shopping requests in English and Hinglish. It queries cross-store inventories, verifies fresh stock, optimizes the shopping basket across multiple vendors, and provides a visual Store Journey Audit proving real net cost savings (demonstrating a ₹267 saving on standard orders after accounting for individual store delivery fees).
    \item \textbf{Reverse RFQ Marketplace and Merchant Quotation Engine (Module B):} Developed a reverse bidding system where customers can post custom or bulk shopping requirements with target budgets and delivery deadlines. Nearby store owners submit itemized line-item quotations detailing unit prices, stock confirmation, and delivery fees. A multi-criteria ranking engine sorts offers dynamically by lowest price, fastest delivery, highest merchant rating, or balanced value, and seamlessly converts accepted quotations into confirmed orders.
    \item \textbf{Hybrid Recommendation System:} Combined collaborative filtering (user-item cosine similarity), content-based vectors (16-dimensional product attributes), and contextual time-of-day scoring to produce personalized product rails on the homepage without relying on heavy external machine learning services.
    \item \textbf{Operational Demand Prediction:} Implemented an automated background cron pipeline combining exponential smoothing and linear regression to forecast hourly product demand and surface restocking alerts on the merchant dashboard.
    \item \textbf{AI Customer Support CRM:} Built a ticketing dashboard featuring automated intent classification, confidence scores, real-time ticket messaging, and sensitive PII data masking to protect user privacy.
    \item \textbf{Cloud Deployment and Verification:} Deployed the full-stack system on Render connected to MongoDB Atlas, verified all functional requirements, and achieved a 100\% pass rate on automated algorithmic test suites.
\end{enumerate}

\section{Practical Impact and Regional Relevance}

In Tier-2 and Tier-3 Indian cities like Imphal, local neighborhood merchants often maintain superior access to fresh agricultural produce and regional grain markets compared to centralized corporate warehouses. However, these local stores typically lack digital infrastructure to reach customers efficiently.

GroDealz demonstrates that an open multi-store model can connect neighborhood stores directly with local consumers. By enabling reverse quotation bidding, local shops can compete transparently on price and service quality, while customers gain access to better prices and customized bulk purchases that standard quick-commerce apps do not accommodate.

\section{Current Limitations}

While the current platform is fully functional as a prototype, several limitations should be noted:
\begin{itemize}
    \item \textbf{Payment Integration:} Financial transactions currently default to Cash on Delivery (COD) or simulated checkout. A commercial deployment requires integration with a live payment gateway (such as Razorpay or UPI Auto-Pay) with escrow capabilities to hold funds during quotation negotiation.
    \item \textbf{Geographic Distance Modeling:} The current system computes Euclidean distances between customer and store coordinates. In a dense urban setting, real-world road network travel times using routing APIs (such as OpenStreetMap or Google Distance Matrix) would provide more precise delivery estimates.
    \item \textbf{Native Mobile Experience:} The current implementation is a responsive web application (SPA). While it performs well on mobile browsers, native Android and iOS applications would offer deeper device integration, such as native push notifications and background geolocation tracking.
\end{itemize}

\section{Future Scope}

Building on the foundation established in this report, several exciting extensions can be pursued:
\begin{enumerate}
    \item \textbf{UPI Escrow Smart Contracts:} Implementing automated digital escrow accounts where a customer's budget is temporarily locked upon posting an RFQ and released to the merchant only after delivery is confirmed via an OTP.
    \item \textbf{Voice-Based Multilingual Assistant:} Expanding the natural language parser to accept voice input in local languages (such as Manipuri / Meiteilon and Hindi) to make grocery ordering accessible to non-English-speaking users and elder citizens.
    \item \textbf{Delivery Rider Pooling and Route Optimization:} Integrating a shared rider dispatch system where a single delivery agent picks up items from two adjacent shops on behalf of the customer, further reducing multi-store delivery costs.
    \item \textbf{Automated Computer Vision Shelf Auditing:} Allowing store owners to snap a photograph of their store shelves to automatically update stock counts using object detection models.
\end{enumerate}
